\documentclass[11pt, a4paper]{article}

\usepackage[utf8]{inputenc}
\usepackage[ngerman]{babel}
\usepackage{amsmath}
\usepackage{amssymb}
\usepackage{amsthm}
\usepackage{booktabs}
\usepackage{hyperref}
\usepackage{geometry}
\geometry{a4paper, margin=2.5cm}

\usepackage{microtype}
\usepackage{listings}
\usepackage{xcolor}
\usepackage{listingsutf8}

\newtheorem{definition}{Definition}
\newtheorem{theorem}{Theorem}
\newtheorem{axiom}{Axiom}
\newtheorem{lemma}{Lemma}

\begin{document}
	\pagenumbering{roman}
	\begin{titlepage}
		\centering
		\vspace*{4cm}
		
		{\scshape\Large Epistemische Formalisierung\par}
		\vspace{1.5cm}
		
		{\huge\bfseries Incoherence is all you need.\par}
		\vspace{1cm}
		
		{\Large\itshape TrigBiPy: Eine formale epistemische Grammatik zur algorithmengestützten Erkenntnisgewinnung\par}
		\vspace{3cm}
		
		{\Large\textbf{Jason A. Schurawel}\par}
		\vfill
		
		{\large \today\par}
	\end{titlepage}	
	\pagenumbering{arabic}
	
	\newpage
	\tableofcontents
	\newpage
	
	\renewcommand{\abstractname}{Abstract}
	\addcontentsline{toc}{section}{A\quad Abstract}
	\begin{abstract}
		\begin{quote}
		Der vorliegende Artikel skizziert die Architektur von „TrigBiPy“, einem formalen System zur Validierung menschlicher Erkenntnisprozesse. Entgegen traditionellen Ansätzen des Personal Knowledge Managements (PKM) und probabilistischer Sprachmodelle postuliert TrigBiPy, dass Erkenntnisgewinn als deterministisches System epistemischer Zustandsübergänge (Morphismen) beschrieben werden kann. Durch die Reduktion auf vier fundamentale Kategorien (Quelle, Beobachtung, Theorie, Diskurs) wird eine geschlossene algebraische Struktur definiert, wobei Erkenntnis sich asymptotisch in Attraktoren stabilisiert. Das System fungiert dabei nicht als semantischer Wahrheitsgenerator, sondern als syntaktischer Compiler (epistemischer Linter), der logische Inkohärenzen aufdeckt und Argumente algorithmisch bewertet. Implementiert als \emph{Knowledge-as-Code}-Architektur lässt sich die Kohärenz durch SAT-Solver, Lean oder Prolog formal validieren.
		\end{quote}
	\end{abstract}
	
	
	\section{Einleitung: Die strukturelle Krise der Wissensrepräsentation}
	Historische Konzepte der "Augmented Intelligence" (etwa Bushs Memex oder Nelsons Xanadu) scheitern in der Praxis oft an der kognitiven Überlastung des Nutzers. Wenn alle Informationen frei miteinander verknüpfbar sind, entsteht ein semantischer Sumpf ohne logische Verifikationsinstanz. Die Suche nach der perfekten Methode für ein spezifisches Forschungsobjekt ist oft ein Holzweg, da sich die Methode nicht am Inhalt des Objekts festmachen lässt. 
	
	Das TrigBiPy-Modell begegnet diesem Problem durch die These, dass Erkenntnisgewinnung kein ungeordneter kreativer Akt ist, sondern eine kontrollierte Dynamik der Inkohärenzreduktion. Das System erzeugt keine autonome Wahrheit, sondern prüft, ob epistemische Konstruktionen innerhalb einer definierten Grammatik zulässig sind. TrigBiPy fungiert somit als epistemischer Linter – vergleichbar mit Theorembeweisern wie \emph{Lean} oder \emph{Coq}, in denen Beweise als Programme und Aussagen als Typen behandelt werden.
	
	\section{Architektur des Modells: Epistemische Ontologie und Syntax}
	
	\subsection{Kategoriale Ontologie und strukturelle Exklusion}
	Das System operiert mit vier fundamentalen epistemischen **Kategorien** ($\mathcal{C} = \{S, O, T, D\}$). Konkrete Instanzen bezeichnen wir mit Kleinbuchstaben: $s \in S$, $o \in O$, $t \in T$, $d \in D$.
	
	\begin{table}[ht]
		\centering
		\begin{tabular}{l c p{10cm}}
			\toprule
			\textbf{Kategorie} & \textbf{Symbol} & \textbf{Beschreibung} \\
			\midrule
			\textbf{Source} & $S$ & Uninterpretiertes Rohmaterial (z. B. Primärtexte, unstrukturierte Daten). Quellen besitzen im System keine inhärente Bedeutung. \\
			\addlinespace
			\textbf{Observation} & $O$ & Eine interpretierte Selektion. Jede Beobachtung ist theoriegeladen; eine interpretationsfreie Wahrnehmung existiert nicht. \\
			\addlinespace
			\textbf{Theory} & $T$ & Ein stabilisiertes Interpretationsmodell. Theorien sind Mengen von strukturierenden Aussagen. \textbf{Axiome ($A$)} bilden keine separate Klasse, sondern sind als unveränderliche Elemente eine Teilmenge der Theorie ($A \subseteq T$, mutable = false). \\
			\addlinespace
			\textbf{Discourse} & $D$ & Der Zustand der expliziten Inkohärenz oder Spannung. Diskurse repräsentieren strukturelle Unvereinbarkeiten. Unlösbare Diskurse bezüglich einer gegebenen Menge $S$ bilden dabei die Teilmenge der \textbf{Unresolvables ($U$)} ($U \subseteq D$). \\
			\bottomrule
		\end{tabular}
		\caption{Die vier fundamentalen epistemischen Kategorien in TrigBiPy.}
	\end{table}
	
	\textbf{Exklusion statischer Forschungsprodukte:} Traditionelle Endpunkte, wie man sie oft aus dem Personal Knowledge Management (PKM) als finale Forschungsprodukte oder Ergebnisse kennt, werden in TrigBiPy vollständig aus der formalen epistemischen Grammatik ausgeschlossen. Das fertige Forschungsprodukt ist in diesem Paradigma streng von der eigentlichen, iterativen Erkenntnisgewinnung zu trennen. Das System verlangt eine rein bipartite Enddynamik zwischen Theorie ($T$) und Diskurs ($D$). Finale Resultate sind lediglich kommunikative Export-Artefakte, die selbst keinen strukturellen Erkenntnisüberschuss mehr innerhalb der Systemalgebra generieren. Sie stellen demnach keine zulässige epistemische Kategorie für morphologische Operationen dar.

	\subsection{Die Invarianten der epistemischen Grammatik}
	\begin{itemize}
		\item \textbf{Gesetz der Unerschaffbarkeit}: Eine Quelle ($s \in S$) kann niemals das Produkt einer Systemoperation sein.
		\item \textbf{Gesetz der blinden Daten}: Quellen ($S$) können nicht direkt miteinander interagieren ($S \times S \to \emptyset$).
		\item \textbf{Gesetz der binären Reduktion und Kommutativität}: Die Reihenfolge der Operanden verändert den epistemischen Endzustand nicht. Alle komplexen Übergänge sind auf primitive binäre Operatoren reduzierbar.
	\end{itemize}
	
	\subsection{Der Abbildungsraum und die Bifurkationsmatrix}
	Wir formalisieren den Möglichkeitsraum der Erkenntnisgewinnung. Sei $\mathcal{C} = \{S, O, T, D\}$ die Menge der epistemischen Kategorien.
	
	\begin{definition}[Epistemischer Möglichkeitsraum]
		Sei $\pi: \mathcal{C} \times \mathcal{C} \to \mathcal{P}(\mathcal{C})$ die Abbildungsfunktion des epistemischen Möglichkeitsraums, welche jedem Paar von Kategorien eine Teilmenge zulässiger Folgekategorien zuordnet.
	\end{definition}
	
	Die zulässigen Projektionen des Möglichkeitsraums lassen sich für die geordnete Basis $(S, O, T, D)$ in folgender symmetrischer Matrixdarstellung $M_\pi$ abbilden:
	
	\[
	M_\pi = \begin{pmatrix}
		\emptyset & \emptyset & \{O\} & \{O\} \\
		\emptyset & \{D, T\} & \{D, T\} & \{D, T\} \\
		\{O\} & \{D, T\} & \{D, T\} & \{D, T\} \\
		\{O\} & \{D, T\} & \{D, T\} & \{D, T\}
	\end{pmatrix}
	\]
	
	Die Matrix $M_\pi$ offenbart das formale Kernstück der TrigBiPy-Architektur: Abseits der initialen theorie- oder diskursgeladenen Datenextraktion ($S \times T \to O$ und $S \times D \to O$) mündet jede Verknüpfung im System zwingend in einer Dualität aus Diskurs ($D$) oder Theorie ($T$). Diese \emph{strukturelle Bifurkation} ist kein Mangel an Systempräzision, sondern die mathematische Modellierung menschlicher Kognition. 
	
	Bei der algebraischen Kollision zweier Informationsknoten $x, y \in \mathcal{C}$ muss der Linter stets zwei formale Endzustände $z \in \pi(x, y)$ vorhalten:
	\begin{itemize}
		\item \textbf{Synthese ($z \in T$)}: Formalisiert als $x \land y \vdash_{A} \top$. Die Informationen verifizieren sich logisch widerspruchsfrei unter den bestehenden Axiomen $A$. Die epistemische Spannung entlädt sich, und das Interpretationsmodell wird stabilisiert.
		\item \textbf{Anomalie ($z \in D$)}: Formalisiert als $x \land y \vdash_{A} \bot$. Die Informationen erzeugen einen Widerspruch und somit eine strukturelle Inkohärenz. Das System generiert eine logische Reibung, die nicht gelöscht werden darf, sondern solange als eigener Knoten ($D$) bestehen bleibt, bis eine Axiom-Modifikation oder ein neuer theoretischer Pfad den Konflikt auflöst.
	\end{itemize}
	Die zwingende Existenz beider Endzustände schützt den Erkenntnisgraphen formell vor degenerierten Systemzuständen: Ein Linter, der algebraisch ausschließlich $T$ zuließe, würde Widersprüche systemisch ignorieren (Dogmatismus). Ein Linter, der ausschließlich $D$ zuließe, verhinderte jegliche Attraktorenbildung (Skeptizismus).
	
	\section{Dynamik der Erkenntnis}

	\subsection{Epistemische Kardinalität}

	Um die Kardinalität einer Verknüpfung zu definieren, ohne den primären epistemischen Formalismus des Systems zu überladen, betrachten wir an dieser Stelle lokal den Raum aller theoretisch potenziell erschaffbaren Ausdrücke.

	\begin{definition}[Raum potenzieller Artefakte]
		Sei $\mathbb{W}_{\text{pot}}$ eine exklusiv für diese Metrik definierte unendliche Menge aller sprachlich oder formallogisch konstruierbaren Artefakte.
	\end{definition}

	\begin{definition}[Lokale Typisierung]
		Wir definieren eine Typisierungsfunktion $\mu_{\text{cat}}: \mathbb{W}_{\text{pot}} \to \mathcal{C}$, die einem potenziellen Artefakt eine der vier Kategorien aus $\mathcal{C} = \{S, O, T, D\}$ zuweist.
	\end{definition}

	Wenn wir nun zwei Artefakte verknüpfen, generiert der menschliche Verstand (oder ein Algorithmus) völlig neue Elemente in $\mathbb{W}_{\text{pot}}$.

	\begin{definition}[Semantischer Generierungsraum]
		Sei $\Gamma_{\text{gen}}: \mathbb{W}_{\text{pot}} \times \mathbb{W}_{\text{pot}} \to \mathcal{P}(\mathbb{W}_{\text{pot}})$ eine generative Funktion, die aus zwei Ursprungsartefakten $a, b \in \mathbb{W}_{\text{pot}}$ eine Menge von inhaltlich plausiblen Folgeartefakten generiert.
	\end{definition}

	Ein generiertes Artefakt $c \in \Gamma_{\text{gen}}(a,b)$ ist zunächst reine Prosa. Es wird erst dann zu einem validen Erkenntnisschritt, wenn die Kategorie, die wir ihm zuweisen ($\mu_{\text{cat}}(c)$), syntaktisch durch die Grammatik ($\pi$) erlaubt ist.

	\begin{definition}[Valide Ableitungsmenge]
		Die lokal definierte Menge der syntaktisch legalen Folgeartefakte $\mathbb{V}_{\text{val}}$ ist:
		\[
		\mathbb{V}_{\text{val}}(a, b) := \{ c \in \Gamma_{\text{gen}}(a, b) \mid \mu_{\text{cat}}(c) \in \pi(\mu_{\text{cat}}(a), \mu_{\text{cat}}(b)) \}
		\]
	\end{definition}

	Jetzt ist die Kardinalität wasserdicht definiert, da sie exakt misst, wie viele völlig neue, gültige Elemente in unsere Graph-Datenbank geschrieben werden könnten.

	\begin{definition}[Epistemische Kardinalität $\kappa$]
		Die Kardinalität einer epistemischen Verknüpfung entspricht der Mächtigkeit der validen Ableitungsmenge:
		\[
		\kappa(a, b) := |\mathbb{V}_{\text{val}}(a, b)|
		\]
	\end{definition}

	\begin{itemize}
		\item \textbf{Totale epistemische Blockade ($\kappa = 0$):}  
		Sei $a \in \mathbb{W}_{\text{pot}}$ die Beobachtung ($\mu_{\text{cat}}(a) = O$) „Die Straße ist nass“ und $b \in \mathbb{W}_{\text{pot}}$ die Theorie ($\mu_{\text{cat}}(b) = T$) „Die Winkelsumme im Dreieck ist 180 Grad“. Die Generierungsfunktion liefert keine logisch plausiblen neuen Sätze: $\Gamma_{\text{gen}}(a,b) = \emptyset$. Folglich ist die Menge der validen Ableitungen leer: $\mathbb{V}_{\text{val}}(a,b) = \emptyset \implies \kappa = 0$. (System befindet sich in Stagnation).
		
		\item \textbf{Eindeutige Synthese ($\kappa = 1$):}  
		Sei $a$ die Beobachtung ($\mu_{\text{cat}}(a) = O$) „Es regnet“ und $b$ die Theorie ($\mu_{\text{cat}}(b) = T$) „Regen macht Straßen nass“. Die Generierungsfunktion erschafft genau einen neuen Satz $c =$ „Die Straße wird nass“. Wir weisen $c$ die Rolle einer Theorieerweiterung zu: $\mu_{\text{cat}}(c) = T$. Da $T \in \pi(O, T)$ erlaubt ist, gilt $\kappa = 1$. (Dies impliziert schnelle System-Konvergenz).
		
		\item \textbf{Endliche Ambiguität ($\kappa > 1$):}  
		Die Funktion $\Gamma_{\text{gen}}$ erschafft aus einer Beobachtung und Theorie zwei widersprüchliche Folgeartefakte $c_1, c_2$ mit $\mu_{\text{cat}}(c_i) = T$. Weil beide in Konflikt stehen, erzwingt das System eine Abzweigung, und $\kappa = 2$. Dies triggert den Diskurs.
		
		\item \textbf{Offene Ambiguität ($\kappa = \infty$):}  
		Unter einer vagen Theorie ohne Restriktionen ist die Menge der plausiblen Sätze $\Gamma_{\text{gen}}(a,b)$ unendlich groß. Das System befindet sich in epistemischer Sättigung ($\kappa = \infty$) und erfordert zwingend axiomatische Einschränkungen.
	\end{itemize}

	Damit ist formal bewiesen, dass TrigBiPy nicht auf einer statischen Datenbank aufbaut. Das System lässt bei jedem Schritt eine Unendlichkeit an Prosa ($\mathbb{W}_{\text{pot}}$) zu, diszipliniert die Wissensbildung jedoch radikal durch den Schnitt aus semantischer Sinnhaftigkeit ($\Gamma_{\text{gen}}$) und strenger Grammatik ($\pi$).
	
	\subsection{Der iterative Zyklus: Charakterisierung von Transitionen und Zuständen}
	Um den iterativen Prozess formal zu fassen, muss zwingend zwischen den abstrakten Kategorien $\mathcal{C} = \{S, O, T, D\}$ und den tatsächlichen, zeitabhängigen Mengen der instanziierten Elemente im Systemgraphen unterschieden werden. Alle im System verarbeiteten Elemente sind Teilmengen des universellen Raums potenzieller Artefakte $\mathbb{W}_{\text{pot}}$.
	
	Wir definieren das globale System zum diskreten Zeitpunkt $t$ als das Tupel der aktuell existierenden, kategorienspezifischen Teilmengen: $E_t = (\mathfrak{S}_t, \mathfrak{O}_t, \mathfrak{T}_t, \mathfrak{D}_t)$, wobei $\mathfrak{S}_t, \mathfrak{O}_t, \mathfrak{T}_t, \mathfrak{D}_t \subset \mathbb{W}_{\text{pot}}$ und $\mathfrak{T}_t \neq \emptyset$ (da mindestens ein Axiom existieren muss). Die Epistemik operiert als zeitdiskretes dynamisches System $E_{t+1} = F(E_t)$.
	Löschen von Knoten ist streng verboten (Wahrung der Historizität). Die Systemmengen wachsen demnach streng monoton: Für jede Kategorie $X \in \mathcal{C}$ und ihre zugehörige Menge $\mathfrak{X}_t$ gilt zwingend $\mathfrak{X}_{t+1} \supseteq \mathfrak{X}_t$. Anstelle einer Entitätslöschung bei Falsifikation generiert das System eine dynamische Angriffsrelation $\mathcal{R}_t \subseteq \mathbb{W}_{\text{pot}} \times \mathbb{W}_{\text{pot}}$ im Sinne eines Argumentationsframeworks nach Dung. Ein ungültiges Element verbleibt stets in seiner Menge, wird aber durch eine Kante $(a, b) \in \mathcal{R}_t$ (Argument $a$ greift Argument $b$ an) epistemisch inaktiviert.
	
	Eine formale Operation $op(x, y) \to z$ projiziert zwei bestehende Elemente $x \in \mathfrak{X}_t, y \in \mathfrak{Y}_t$ (mit korrespondierenden Kategorien $X, Y \in \mathcal{C}$) auf ein Zielelement $z \in \mathbb{W}_{\text{pot}}$. Die Typisierung von $z$ wird durch die syntaktische Matrix auf eine Zielkategorie $Z \in \pi(X, Y)$ abgebildet ($\mu_{\text{cat}}(z) = Z$). Die dem Ziel $Z$ zugewiesene Menge definieren wir als $\mathfrak{Z}_t$. Die topologische Dynamik unterscheidet exakt zwischen zwei Ziel-Szenarien für die resultierende Menge $\mathfrak{Z}_{t+1}$:
	\begin{itemize}
		\item \textbf{Verknüpfung mit bestehendem Element:} Das Ziel ist ein bereits im Graphen instanziiertes Element $z \in \mathfrak{Z}_t$. Die Graphenmengen bleiben kardinal identisch ($\mathfrak{Z}_{t+1} = \mathfrak{Z}_t$), es wird lediglich eine neue topologische Kante gezogen.
		\item \textbf{Erschaffung eines neuen Elements:} Die Operation generiert ein völlig neues Artefakt $z_{new} \in \mathbb{W}_{\text{pot}}$ mit $z_{new} \notin \mathfrak{Z}_t$. Das System integriert dieses Element, woraus ein explizites Mengenwachstum resultiert: $\mathfrak{Z}_{t+1} = \mathfrak{Z}_t \cup \{z_{new}\}$.
	\end{itemize}
	\begin{definition}[Durchführung der Operationen]
	Die \emph{epistemische Entropie} $H_E$ des Systems wird als monoton steigende Funktion über die unlösbaren Inkohärenzen im Graphen definiert, $H_E \propto |\mathfrak{D}_t|$. Das System operiert deterministisch, indem es systematisch prüft, welche Operationen syntaktisch in $\mathcal{C}$ möglich sind. Aus dieser Menge werden iterativ jene Operationen isoliert, die semantisch sinnvoll sind. Diese validierten Operationen ($\mathcal{O}_{val}$) werden anschließend algorithmisch angewandt, mit dem expliziten Ziel der Entropie- und Inkohärenzminimierung. Mathematisch lässt sich dies als iterativer Optimierungsprozess fassen:
	\[
	op^* = \arg\min_{op \in \mathcal{O}_{val}} \Delta H_E(E_t \cup \{op\})
	\]
	Durch diese zielgerichtete Reorganisation erzwingt das System die Synthese neuer kohärenter Theorie-Elemente ($t_{new}$), wodurch die verbleibende Diskursmenge strukturell aufgelöst und die globale Entropie minimiert wird ($\frac{dH_E}{d\mathfrak{T}} < 0$).
	\end{definition}
	
	\subsection{Wissenschaftstheoretische Implikationen}
	\begin{itemize}
		\item \textbf{Epistemische Konvergenz und Attraktorstruktur:}
		Wissenschaftliche Stabilisierung ist kein monotoner Akkumulationsprozess von Daten, sondern ein rekursiver Attraktorprozess im Raum der epistemischen Zustände. Unter der Annahme einer begrenzten Injektion fundamental neuer Quellen ($|S_{\text{neu}}| < \infty$) und einer endlichen Diskursgenerierung konvergiert das dynamische System:
		\[
		\lim_{t \to \infty} D_t \to \emptyset \implies \exists T^* : F(T^*) = T^*
		\]
		Theorien $T^*$ fungieren als \emph{Fixpunkte} (Attraktoren) im dynamischen Raum der Erkenntnis. Wissenschaftliche Stabilisierung ist formal die strukturelle Kompression von Beobachtungen und Diskursen in ein asymptotisches Interpretationsmodell $T^*$.
		
		\item \textbf{Korrektur der empiristischen Induktion:}
		Diese mathematische Formulierung korrigiert eine klassische empiristische Fehlannahme der Wissenschaftstheorie: Neue Theorien ($t_{new} \in \mathfrak{T}_{t+1}$) entstehen nicht durch naive, direkte Induktion aus nackten Daten ($op(o_1, o_2) \to t_{new}$ ist strukturell unzureichend für einen Paradigmenwechsel), sondern zwingend durch die Auflösung formaler Spannungen. Wie historische Paradigmenwechsel (etwa der Übergang zur Relativitätstheorie) verdeutlichen, führt eine anomale Beobachtung $o_{new} \in \mathfrak{O}_t$ nicht isoliert zu einem neuen Paradigma. Vielmehr kollidiert $o_{new}$ zwingend mit dem bestehenden Modell $t_{alt} \in \mathfrak{T}_t$ und erzeugt durch die Verknüpfung formal einen Diskursknoten $op(o_{new}, t_{alt}) \to d_{new} \in \mathfrak{D}_{t+1}$. Erst aus diesem spannungsgeladenen Diskursraum geht unter den strukturellen Limits der Axiomatik ($a \in A \subseteq \mathfrak{T}$) die neue Theorie hervor ($op(d_{new}, a) \to t_{new}$). Das Element $t_{new}$ etabliert anschließend einen logischen Angriff $(t_{new}, t_{alt}) \in \mathcal{R}_{t+n}$.
		
		\item \textbf{Inkohärenz als primäre Ressource:}
		Erkenntnis ist somit keine Wahrscheinlichkeitsmodellierung (wie es statistische LLMs simulieren) und auch keine reine Suchabfrage in Wissensgraphen (PKM). Es ist ein \emph{constraint-basiertes Rewriting-System}, dessen primäre Berechnungsressource die Inkohärenz ist. Inkohärenz ($D$) ist kein "Fehler" im System, sondern das essenzielle Signal, das neue theoretische Stabilisierungen überhaupt erst erzwingt. Das TrigBiPy-System optimiert folglich nicht primär auf unmessbare absolute Wahrheiten, sondern minimiert als formale Zustandsmaschine iterativ epistemische Inkonsistenzen.
	\end{itemize}
	
	\subsection{Diskurs-Sättigung und Folgen}
	
	Wir greifen die eingangs definierte Teilmenge der unlösbaren Diskurse auf, formalisiert als $\mathfrak{U}_t \subseteq \mathfrak{D}_t$. Es stellt sich die systemische Frage, ob das dynamische System $F(E_t)$ bei einer beliebigen, aber fixierten Quellenmenge $\mathfrak{S}_t = const$ zwingend in einen vollständig kohärenten Attraktorraum konvergiert, sodass $\lim_{k \to \infty} \mathfrak{U}_{t+k} = \emptyset$.
	
	\begin{theorem}[Epistemische Stagnation]
		Für eine statische Quellenmenge $\mathfrak{S}_t \neq \emptyset$ existiert nicht notwendigerweise ein Fixpunkt mit $\mathfrak{U}^* = \emptyset$. Wenn für die Menge aller zulässigen Folgeoperationen über $E_t$ gilt, dass kein neues theoriebildendes Element $t_{new} \in \mathfrak{T}_{t+1}$ synthetisiert werden kann, welches einen validen Angriff $(t_{new}, u) \in \mathcal{R}_{t+1}$ auf die unlösbaren Diskurse generiert, konvergiert das System in eine Stagnation. Die Teilmenge der unlösbaren Diskurse wird invariant: $\mathfrak{U}_{t+1} = \mathfrak{U}_t \neq \emptyset$.
	\end{theorem}
	
	In der realen Wissenschaftsgeschichte manifestiert sich dieser Zustand als permanente Koexistenz unvereinbarer Theorien, die empirisch unterdeterminiert sind. Verschiedene Elemente $t_1, t_2 \in \mathfrak{T}_t$ leiten sich widerspruchsfrei aus derselben Beobachtungsmenge $\mathfrak{O}_t$ ab. Da jedoch keines der Elemente stark genug ist, um das andere erfolgreich anzugreifen ($(t_1, t_2) \notin \mathcal{R}_t$ und $(t_2, t_1) \notin \mathcal{R}_t$), erzeugen sie in ihrer gegenseitigen Verknüpfung zwingend ein unlösbares Knotenelement $op(t_1, t_2) \to u \in \mathfrak{U}_t$. Ein prominentes physikalisches Beispiel ist die Quantenmechanik: Verschiedene Interpretationen (etwa Kopenhagener Deutung vs. Viele-Welten-Theorie) operieren erfolgreich auf der exakt gleichen Quellbasis $\mathfrak{S}_t$, verbleiben jedoch systemisch in einem unauflösbaren Spannungsverhältnis.
	
	\begin{theorem}[Symmetriebrechung und kaskadierender Kollaps]
		Sei $E_t$ ein System in Sättigung ($\mathfrak{U}_t = const > 0$). Wird nun ein fundamentales Element in die Quellenmenge injiziert, sodass $\mathfrak{S}_{t+1} = \mathfrak{S}_t \cup \{s_{new}\}$, resultiert daraus iterativ eine neue Beobachtung $o_{new} \in \mathfrak{O}_{t+2}$. Kollidiert $o_{new}$ asymmetrisch mit den bestehenden Theorien (z. B. $op(o_{new}, t_1) \to d_{new}$ und $op(o_{new}, t_2) \to t_{new}$), bricht die formale Symmetrie der Unterdeterminierung. Das neue Element $t_{new}$ formt nun einen validen Angriff $(t_{new}, t_1) \in \mathcal{R}_{t+n}$.
	\end{theorem}
	
	Die Injektion von $s_{new}$ wirkt als Katalysator: Zuvor inkompatible Elemente aus $\mathfrak{U}_t$ werden durch die asymmetrische Schwächung von $t_1$ (Invalidierung durch die Angriffsrelation) zugunsten von $t_2$ abrupt auflösbar. Diese plötzliche algebraische Resolution einer zuvor saturierten Diskursmenge beschreibt formell den Prozess, der wissenschaftshistorisch als revolutionärer Paradigmenwechsel bekannt ist: Ein schlagartiger topologischer Graphenkollaps und der plötzliche Wechsel ganzer Theorienkonstrukte.
	
	\subsection{Kategorientheoretische Formalisierung und Attraktoren}
	Um die strukturellen Eigenschaften der Wissensdynamik ($\Delta H_E$) formell exakt zu fassen, modellieren wir TrigBiPy als streng typisierte Kategorie $\mathcal{E}$, innerhalb derer Erkenntnis als iterativer Konvergenzprozess durch epistemische Endofunktoren abgebildet wird. Dies erlaubt es uns, abstrakte Transformationen rigoros nachzuweisen und semantische Konsistenz strukturell zu erzwingen.
	
	\textbf{Definition der Kategorie $\mathcal{E}$:}
	\begin{itemize}
		\item \textbf{Objekte ($Ob(\mathcal{E})$):} Die fundamentalen Kategorien $\mathcal{C} = \{S, O, T, D\}$ fungieren nicht bloß als lose Datentypen, sondern als formale Objekte im kategorientheoretischen Sinne. Eine spezifische Instanziierung des Systemgraphen zu einem Zeitpunkt $t$ ist ein konkretes Morphismennetzwerk zwischen diesen Objekten.
		\item \textbf{Morphismen ($Hom(X, Y)$):} Jeder bewusste Erkenntnisschritt (z. B. eine Beobachtung aus einer Quelle generieren, oder zwei Theorien in einen Diskurs überführen) entspricht einem Morphismus. Die epistemische Matrix $M_\pi$ definiert dabei die Menge aller legalen strukturellen Übergänge in $\mathcal{E}$.
	\end{itemize}
	
	\textbf{Der Diskurs als Pullback-Obstruction:}
	Die formelle Eleganz der kategorientheoretischen Betrachtung offenbart sich in der Definition der Diskursknoten ($D$). Wenn zwei getrennte theoriegeladene Beobachtungen $o_1$ und $o_2$ im Erkenntnisprozess nicht widerspruchsfrei unter einem gemeinsamen theoriebildenden Pullback-Objekt $t$ vereint werden können, kommutiert das entsprechende Diagramm nicht. 
	
	In traditionellen Systemen (wie simplen Wissensgraphen) würde diese Inkompatibilität schlicht ignoriert oder gelöscht. In TrigBiPy jedoch materialisiert sich das Scheitern der Kommutativität (die "Pullback Obstruction") als eigenständiges formales Objekt: der Diskurs $D$. Ein Diskurs ist somit keine Fehlerstelle in der Datenbank, sondern das strukturelle Maß eines Nicht-Null-Wertes der \emph{epistemischen Krümmung}. Er zwingt das System formell dazu, den globalen Graphen solange algebraisch umzuformen (durch die Synthese neuer, höherer Theorien), bis die Krümmung aufgelöst und Kommutativität wiederhergestellt ist.

	\section{Angewandte Implementierung und Knowledge-as-Code}
	
	\lstset{
		language=Python,
		basicstyle=\ttfamily\footnotesize,
		keywordstyle=\color{blue},
		stringstyle=\color{red},
		commentstyle=\color{green!60!black},
		breaklines=true,
		frame=single,
		extendedchars=true,
		inputencoding=utf8/latin1
	}

	\subsection{Proof of Concept: Minimum Viable Product}
	
	\lstinputlisting{trigbipy_mvp.py}

	\subsection{Algorithmische Validierung und Knowledge-as-Code}
	Das formale Gerüst von TrigBiPy erlaubt eine direkte algorithmische Implementierung nach dem \emph{Knowledge-as-Code}-Paradigma. Der Graph wird als lokales Dateisystem realisiert (Ordner für Sources, Observations, Theories und Discourses, die als Markdown-Dateien vorliegen). Git dient hierbei als native Infrastruktur: Eine Theorie entspricht einem Branch, ein Diskurs äußert sich als Merge Conflict, und Falsifikation ist ein Revert.
	
	Da das System rein strukturell agiert, lässt sich die epistemische Konsistenz deterministisch überprüfen. Die konkrete Validierung der Architektur ist derzeit strukturell algorithmisch möglich:
	\begin{itemize}
		\item \textbf{Domänenspezifische Algorithmen:} Diese beruhen auf impliziten Axiomen. Bspw. können Simulationen die Sinnhaftigkeit von Operationen prüfen.
		\item \textbf{SAT-Solver:} Prüfung auf logische Widersprüche innerhalb der Constraints (Konsistenz von $T \equiv SAT(C_T)$).
		\item \textbf{Logic Programming (Prolog):} Resolution der Graphenkanten und Inferenzketten.
		\item \textbf{Dependent Type Systems (Lean):} Eine Theorie fungiert als Typ, eine Beobachtung als Term und ein Diskurs als Type Mismatch.
		\item \textbf{LLM-Ansätze:} Large Language Models werden rein heuristisch zur Überprüfung der inhaltlichen Kohärenz des Fließtextes verwendet, ohne jedoch epistemische Autorität zu besitzen.
	\end{itemize}
	
	\subsection{TrigBiPy als Framework für AGI}
	Das aktuelle Paradigma der Künstlichen Intelligenz (primär LLMs) basiert auf der probabilistischen Vorhersage des nächsten Tokens. Diese Systeme scheitern an echter Erkenntnisgewinnung.
	
	TrigBiPy bietet hier einen völlig neuen Lösungsansatz. Kleine, lokale LLMs dienen lediglich als Übersetzer und heuristische Vorschlagsgeneratoren für die Operationen der epistemischen Grammatik. Der Compiler prüft Dung-Angriffsrelationen, Kardinalität und Typensicherheit.
	
	\textbf{Incoherence is all you need.} Inkohärenz ist der Motor der Erkenntnis. Git dient als Infrastruktur: Hypothesen als Branches, Diskurse als Merge Conflicts, und Erkenntnisiterationen als deterministische Commits in der Knowledge-as-Code-Architektur.
	
	In einer hybriden AGI-Architektur übernimmt das neuronale Netz die Intuition (Vorschläge für Kanten), während der TrigBiPy-Graph die harte Logik erzwingt. Halluzinationen werden zu produktiven Diskursen, die das System zur Auflösung zwingt.
	
\end{document}